\section{Composition through marked intermediate terms}
\label{app:composition}

This appendix proves composition for the relation actually defined in
Section~\ref{sec:comparison}. Auxiliary comparisons may have marked
targets; they will be eliminated when the final target is ordinary.
They are not used as a well-founded relation. In particular, an
already-marked term is strictly comparable to itself in this auxiliary
relation, so its strict part must not be substituted for the numbered
relation of Section~\ref{sec:upper}.

\subsection{The auxiliary rules}

A \emph{prefix copy} of
$P_\tau=f^\star_\tau(\vec b;\vec e)$ is
$f^\star_\rho(\vec b;\vec e_{\le j})$, where
$0\le j\le|\vec e|$ and $\rho$ is its required output type.
The prefix is literal, including order and types. Write $P+\vec a$
for the same marked owner with $\vec a$ appended to its extras.

The least finite-derivation pair $\succeq_e,\succ_e$ has all eight
rules of Section~\ref{sec:comparison}, with the subscript $e$ in
every premise, and the same persistent strict rule. It also has
the following seven weak rules. The base matching in \textsc{In}
allows a permutation; its extra matching does not.
\begin{enumerate}
\item[\textsc{In}]
From $f\approx g$, a weak matching of the base occurrences, and
an ordered, equal-typed weak matching of the extras, infer
\[
 \EW{f^\star_\tau(\vec b;\vec e)}
    {g^\star_\tau(\vec c;\vec d)}.
\]
\item[\textsc{Marked1}]
If $f>g$ and $P$, retyped as necessary, weakly dominates every
base and extra argument of $g^\star_\tau(\vec c;\vec d)$,
infer $\EW{P_\tau}{g^\star_\tau(\vec c;\vec d)}$.
\item[\textsc{Marked2}]
Use the same conclusion and whole-source premises when
$f\approx g$ and
$\vec b>_{\mathrm{mul}}^{\succeq_e,\succ_e}\vec c$.
\item[\textsc{Drop}]
From $\EW P t$ infer $\EW{P+\vec a}t$, for any typed extra list.
\item[\textsc{Perm}]
The source extras may be permuted, with each type attached to
its term.
\item[\textsc{Copies}]
If every entry of $\vec p$ is a prefix copy of $P$, then
$\EW{P+\vec p}t$ implies $\EW P t$.
\item[\textsc{PSelect}]
The \textsc{Select} rule may use a tuple of prefix copies of its
owner in place of the tuple of whole-owner copies. The selected
term remains an actual base or extra argument, and its full
original-spine call must succeed.
\end{enumerate}
Every original derivation embeds in this pair. Structural
induction using \textsc{Var}, \textsc{Abs}, \textsc{Fun}, and
\textsc{In} proves $\EW q q$ for every marked or ordinary term.

\subsection{A syntactic refinement relation}

We give explicitly the syntax relation needed to eliminate the
extra rules. It is useful because a staged special application
can expose one of the copies substituted at an earlier stage.
This situation is a syntactic refinement of the owner even
when it is not one full application of the initially selected
term.

Define $q\sharpens s$ by the following mutually inductive clauses.
They are finite trees, not a greatest fixed point.
Variables refine only themselves; lambdas refine componentwise;
ordinary heads have exactly the same symbol and positionwise
refining base lists. At a marked head,
\[
 f^\star_\tau(\vec b;\vec e)
       \sharpens f^\star_\tau(\vec c;\vec d)
 \tag{S}
\]
holds if $\vec b\sharpens\vec c$ positionwise and every typed
occurrence $d$ in $\vec d$ has one of two finite witnesses:
\begin{enumerate}[label=(\roman*)]
\item an actual occurrence $e$ of the same type in $\vec e$
with $e\sharpens d$;
\item $f^\star_{\operatorname{type}(d)}(\vec b;\vec e)
\sharpens d$.
\end{enumerate}
The second alternative is a premise tree of smaller height in
this mutual definition. No order or multiplicity matching of
extras is required. Head equivalence is insufficient in (S):
the symbol itself is preserved.

\begin{lemma}[Refinement syntax]\label{lem:sharp-syntax}
The relation $\sharpens$ is reflexive and transitive. It commutes
with renaming, opening, and simultaneous substitution by
pointwise refining environments. It preserves retyping of a
marked root and reflects a successful partial mark:
\[
 q\sharpens s,\quad \mk(s)=s'
 \quad\Longrightarrow\quad
 \mk(q)=q'\text{ and }q'\sharpens s'
 \quad\text{for some }q'.
 \tag{M}
\]
Suppose $q\sharpens s$ and $a_i\sharpens b_i$ in corresponding
types. If $\ap s{\vec b}$ succeeds, then $\ap q{\vec a}$ succeeds
and its result refines the result on the right.
\end{lemma}
\begin{proof}
Reflexivity is induction on the term; a marked extra uses its
own occurrence and the induction hypothesis. Transitivity is
mutual induction on the right refinement tree, simultaneously
for bases and the two kinds of extra witness. In the existing
extra case, invert the matching extra of the first refinement.
If it was supplied by an existing occurrence, compose the two
smaller refinements. If it was supplied by the whole owner,
compose that smaller witness with the right extra witness.
In the whole-owner case, apply the marked induction hypothesis
at the requested retyping. The resulting extra witness again
has exactly one of the two forms (i)--(ii). The variable,
lambda, and ordinary-head cases just compose their components.
This proves transitivity without any comparison judgment.

For two-environment substitution, induct over the same mutual
syntax tree. A variable uses its environment witness; bound
variables use a common fresh variable. Ordinary heads and
marked bases retain their positions. An existing extra remains
an existing substituted occurrence. A whole-owner witness is
the corresponding smaller refinement after substitution; the
substituted owner is precisely the owner of the substituted
extra list. Renaming and opening are instances of this proof.
Retyping leaves every base and extra premise unchanged.

For (M), refinement has the same root constructors on both
sides. Descend through their common lambdas. An ordinary head
is changed into the same marked head on both sides, with empty
extras. Marked heads are unchanged by both marks. A variable
on the right cannot have a successful mark.

For full application, follow the common original lambda
spines. If the tuple stops within the spine, the two-environment
substitution result applies. If it proceeds to a head, success
on the right excludes a variable head. Corresponding ordinary
or marked heads therefore remain, and the unconsumed arguments
are appended on both sides. Their positionwise refinements
supply existing-extra witnesses. Old existing-extra witnesses
remain valid after appending. In an old whole-owner witness,
the enlarged source refines the old source by retaining only
the old extras; compose this with that witness. All substitutions
within old bases and extras have already been covered.
\end{proof}

Several useful instances deserve to be stated. A marked owner
refines itself with any extras omitted or permuted. It refines
itself with any finite list of its prefix copies appended:
the old extras use (i), and each appended prefix uses (ii)
followed by the preceding omission instance. Conversely, an
owner with more extras refines the original owner. These
statements allow arbitrary typed multiplicities.

\begin{lemma}[Two-stage prefix calls]\label{lem:prefix-escape}
Let $\vec a,\vec b$ be lists of prefix copies of one marked
owner $P$. Suppose
\[
 \ap q{\vec a}=r,\qquad \ap r{\vec b}=w.
\]
Then either $\ap q{\vec a,\vec b}=w$, or $P\sharpens w$,
at the output type of $w$.
\end{lemma}
\begin{proof}
Follow the original initial lambda spine of $q$. When the
combined tuple does not pass its end, ordinary simultaneous
substitution composition gives the asserted equality. When
it passes that end and the original residual head is
functional, all remaining arguments are appended to that
same head, again giving the equality.

The only remaining case is that the combined call exposes an
original variable. A free variable cannot account for two
successful stages: a still-unconsumed argument at such a
variable would block the second stage. A bound variable can
have received an argument from the first batch. Its value is
a prefix copy of $P$, hence a marked functional owner.
The second call can act on that value. Its result is the
chosen prefix copy, retyped, with a suffix of the later
prefix-copy arguments appended. Every retained old extra is
an existing extra of $P$; every appended copy is refined by a
whole $P$. Clause (S) therefore proves $P\sharpens w$.
This is precisely the case in which full and iterated unary
application would disagree.
\end{proof}

\subsection{Eliminating the auxiliary rules}

\begin{lemma}[Ordinary-target elimination]\label{lem:conservative}
If $q\sharpens s$, $\EW s t$, and $t$ is entirely ordinary,
then $\W q t$. The corresponding statements hold for the
strict judgment, weak matching, strict multiset replacement,
and replacement suppliers.
\end{lemma}
\begin{proof}
Use mutual induction over the five finite comparison and
list-judgment trees. Keep the refinement witness arbitrary.
In \textsc{Var} and \textsc{Abs}, inversion of refinement
gives the same variable or the two body witnesses.
\textsc{Fun} uses the component induction hypotheses.
\textsc{Put} uses (M). For \textsc{Copy}, retype the whole
source witness and apply the induction hypothesis to each
target argument; the head is unchanged. \textsc{Stat} also
uses the multiset induction hypothesis. Opening commutes
with refinement, so \textsc{FAbs} follows from its body
induction hypothesis and the same rule.

For either selection rule, let $a$ be the selected argument
of $s$. If it is a base, refinement supplies an actual
corresponding base $b$ of $q$ with $b\sharpens a$. If it is
an extra, its witness either supplies such an actual extra
$b$ or says that the retyped whole $q$ refines $a$.
Every whole $q$ copy refines the corresponding prefix copy
of $s$: use $q\sharpens s$, then omit the unwanted extras.
In the actual-argument case, Lemma~\ref{lem:sharp-syntax}
gives a successful full call on $b$ with whole $q$ copies
whose result refines the given selected result. Apply the
continuation induction hypothesis and the original
\textsc{Select} rule.
In the whole-owner case, the same application lemma gives
an owner with whole copies appended, refining the selected
result. The original owner refines this expanded owner by
the instances following Lemma~\ref{lem:sharp-syntax}.
Transitivity of refinement therefore lets the continuation
induction hypothesis finish directly.

In \textsc{Drop}, the enlarged owner refines the owner in
the premise by omitting its appended extras. In \textsc{Perm}
it refines the owner with the premise order. In \textsc{Copies},
the unexpanded owner refines the expanded one by the prefix
copy instance. Compose the supplied refinement with this
syntactic witness and use the smaller premise induction
hypothesis.
The conclusions of \textsc{In}, \textsc{Marked1}, and
\textsc{Marked2} are marked and so cannot occur at this
ordinary-target weak judgment.

For a strict judgment, (M) gives the successful mark on the
refining source, and its weak premise is smaller. In the
list judgments, keep the very same partitions and
occurrence maps, apply the weak or strict induction
hypothesis at each occurrence, and retain nonemptiness of
the removed list. These cases complete the mutual induction.
\end{proof}

\begin{lemma}[Compression of extras]\label{lem:extra-compression}
Let $P=f^\star(\vec b;\vec e)$ and
$T=g^\star(\vec c;\vec d)$ have the same output type.
Suppose $f\approx g$, the base lists have an auxiliary weak
matching, and $\EW{P_{\operatorname{type}(d)}}d$ for every
target extra $d$. Then $\EW P T$.
\end{lemma}
\begin{proof}
Let $\vec p$ have, in the order of $\vec d$, one whole
$P_{\operatorname{type}(d)}$ for each occurrence $d$.
\textsc{In} first compares
$f^\star(\vec b;\vec p)$ with $T$.
\textsc{Drop} appends $\vec e$ to this source;
\textsc{Perm} moves $\vec e$ in front of $\vec p$.
Every entry of $\vec p$ is a full-prefix copy of
$f^\star(\vec b;\vec e)$, so \textsc{Copies} discharges
them and gives the assertion. No comparison composition
is used.
\end{proof}

\begin{lemma}[Refinement admissibility]\label{lem:sharp-admissible}
For arbitrary targets,
\[
 q\sharpens s,\quad \EW s t\quad\Longrightarrow\quad\EW q t.
 \tag{A}
\]
The strict and three list-judgment analogues also hold.
\end{lemma}
\begin{proof}
Repeat the mutual derivation induction of
Lemma~\ref{lem:conservative}, returning auxiliary judgments.
All its original-rule, macro, selection, and strict
constructions remain valid with a marked final target.
There are now three additional root cases.

For \textsc{In}, the source refinement preserves the
head and refines each base, so the base matching follows
from the induction hypotheses. A source extra matched
with a target extra $d$ has either an actual-extra
refinement witness $e_q\sharpens e_s$, or a whole-$q$
witness. In the first case the smaller matching
induction hypothesis gives $\EW{e_q}d$; zero-argument
\textsc{Select} then gives $\EW q d$. In the second case
the smaller hypothesis already gives $\EW q d$.
Lemma~\ref{lem:extra-compression} reconstructs the full
marked target with precisely its original extras.

For \textsc{Marked1}, transport each whole-source
premise by the weak induction hypothesis, at its stated
type, and reapply the rule. For \textsc{Marked2}, do
the same and transport its strict multiset by the list
induction hypothesis. This exhausts the extra cases.
Persistence matters in the strict induction: a marked
source is returned by its mark, so (M) applies there
without requiring a fresh syntactic mark.
\end{proof}

\begin{lemma}[Marking transport]\label{lem:extended-marking}
For all auxiliary terms,
\[
 \EW s t\Longrightarrow\EW{\headmark(s)}{\headmark(t)},\qquad
 \EW{\headmark(s)}t\Longrightarrow\EW s t.
 \tag{T}
\]
If $\EW s t$ and $\mk(t)$ is defined, then $\mk(s)$ is defined.
\end{lemma}
\begin{proof}
For the second implication descend through a common
initial lambda spine. At an ordinary head use
\textsc{Put}; at a marked head or a variable spine
the total mark is unchanged. A comparison whose source
is a lambda can use only \textsc{Abs}; a comparison
whose source is a variable is its \textsc{Var}
comparison. This inversion justifies the descent and
also proves reflection of markability.

Prove the first implication by mutual derivation
induction. \textsc{Fun} becomes \textsc{In} with empty
extras. \textsc{Copy} and \textsc{Stat} become,
respectively, \textsc{Marked1} and \textsc{Marked2},
with the same premises. Already-marked target rules
retain their conclusions. \textsc{Put} uses
idempotence. For either selection rule, its smaller
continuation gives the marked selected result above
the marked target. The second implication of (T)
removes the mark from that selected result, so the
same selection applies. \textsc{Abs} and
\textsc{FAbs} use the body hypotheses; the extra-list
macros lift their premise hypotheses. The strict and
multiset cases use the same occurrence witnesses
and $(\headmark)^2=\headmark$.
\end{proof}

\subsection{Covariance of actual full applications}

\begin{lemma}[Two environments]\label{lem:two-environments}
If $\EW s t$ and type-preserving environments $\gamma,\delta$
satisfy $\EW{\gamma(x)}{\delta(x)}$ for every variable,
then $\EW{s\gamma}{t\delta}$. The strict and list-judgment
versions hold as well.
\end{lemma}
\begin{proof}
Induct simultaneously over those judgments, with the
environments arbitrary. A variable uses the supplied
comparison; a binder extends both environments by
the same fresh variable. Head rules preserve all
their finite occurrence maps and use the corresponding
induction hypotheses.

In a selection rule the successful source application,
after $\gamma$, is exactly the original application
result after $\gamma$. This is the one-way
successful-application substitution identity from
Definition~\ref{def:application}. Whole and prefix
copies become literal copies of the substituted owner.
The continuation induction hypothesis has target
$t\delta$, so the same selection rule applies.
\textsc{FAbs} uses lifted environments. The extra-list
macros retain their list and prefix witnesses under
$\gamma$. Finally, a successful mark commutes with
substitution, including its identity case on a marked
root. It supplies the strict induction step. Equal
substituted terms do not collapse indexed occurrences.
\end{proof}

\begin{lemma}[Application covariance]\label{lem:application-covariance}
Suppose $\EW q s$ and $\EW{a_i}{b_i}$ in the respective
argument types. If $\ap s{\vec b}$ succeeds, then
$\ap q{\vec a}$ succeeds and
\[
 \EW{\ap q{\vec a}}{\ap s{\vec b}}.
 \tag{C}
\]
\end{lemma}
\begin{proof}
First check the domains, before decomposing a full call.
Failure is possible precisely when the tuple passes
the original lambda spine and exposes an original
variable. If a source has this form, inversion of its
auxiliary comparison forces the identical lambda
spine and variable on the target: only \textsc{Abs}
and then \textsc{Var} are available. Thus right-hand
full-call success implies left-hand full-call success,
independently of the actual well-typed argument values.

We next prove the one-argument assertion by induction
on $\EW q s$. The \textsc{Var} case is impossible.
For \textsc{Abs}, apply
Lemma~\ref{lem:two-environments} to the body with
$a,b$ assigned to the consumed variable.
\textsc{Fun} and \textsc{In} append $a,b$ to their
respective extra lists and use \textsc{In}.
\textsc{Put} has the same applied source as its
smaller marked-source premise.
\textsc{FAbs} substitutes $a,b$ for the fresh
variable in its premise, again by the two-environment
lemma.

In \textsc{Copy} or \textsc{Marked1}, the left
application appends $a$. The old whole-source
premises survive by \textsc{Drop}, and the new
target extra $b$ is dominated by selecting the new
extra $a$ and using $\EW a b$. \textsc{Marked1}
reconstructs the applied target. For \textsc{Stat}
or \textsc{Marked2}, use \textsc{Marked2} with
the unchanged base multiset. In \textsc{Drop}
and \textsc{Perm}, apply the induction hypothesis
and move the new argument past the old finite extra
list using \textsc{Perm}. In \textsc{Copies}, apply
the hypothesis to the expanded owner, move its
copies after the new argument, and discharge them:
each old prefix is still a prefix of the owner
with the new argument appended.

It remains to check selection. Let $P$ be its owner,
$v$ the selected argument, $\vec p$ its tuple of
prefix copies, and
$\ap v{\vec p}=r$ its result, with $\EW r s$.
Put $Q=P+a$, retyped as necessary. Selecting the
new extra gives $\EW Q b$. Apply the induction
hypothesis to the smaller continuation $\EW r s$,
using the whole $Q$ as left argument and $b$ on
the right. It gives
$\ap r Q=r'$ and $\EW{r'}{\ap s b}$.
Both $\vec p$ and the final whole $Q$ are prefix
copies of $Q$. Lemma~\ref{lem:prefix-escape} gives
either a successful full call $\ap v{\vec p,Q}=r'$
or $Q\sharpens r'$.
In the first case \textsc{PSelect} for the actual
argument $v$ of $Q$ gives the desired comparison.
In the second case
Lemma~\ref{lem:sharp-admissible} gives it directly.
This step uses no comparison cut.

Finally, a successful full call agrees with its
sequence of unary calls. Within the original
lambda spine this is ordinary substitution
composition; beyond a functional head it is
concatenation of extras. The blocked-variable
case has already been excluded for both full
calls. Repeatedly apply the unary result and
use these two identities to obtain (C).
\end{proof}

\subsection{Target transformations and the cut induction}

\begin{lemma}[Target extras]\label{lem:target-macros}
If $\EW s P$ for a marked target $P$, the conclusion
remains true when a suffix of its extras is removed,
when its extras are permuted, or when a finite list
of prefix copies of $P$ is appended.
\end{lemma}
\begin{proof}
For each assertion induct on the comparison derivation.
Only \textsc{In}, \textsc{Marked1}, and
\textsc{Marked2} create the marked target;
the other possible root rules retain their target
and lift the shorter induction conclusion.

For suffix removal in \textsc{In}, remove the
matching source suffix, apply \textsc{In}, and
restore that source suffix by \textsc{Drop}.
For a permutation, permute the matching source
extras, apply \textsc{In}, and restore the source
order by \textsc{Perm}. In either marked precedence
rule, keep just the needed target-extra premises
or reorder their indexing.

For copy addition in \textsc{In}, append to the
source the corresponding source-prefix copies.
Each such source prefix compares with the
requested target prefix by \textsc{In}, using
the restricted old matching. The complete
expanded lists therefore compare by
\textsc{In}; \textsc{Copies} discharges the
added source copies.
For \textsc{Marked1} or \textsc{Marked2}, the
same root rule, with a truncated target extra
list, compares the whole source with each
requested target prefix. These are exactly the
new whole-source premises needed to append
the target copies. No transitivity is invoked.
\end{proof}

\begin{theorem}[Ordinary-final cut]\label{thm:cut}
If $\EW s t$, $\EW t u$, and $u$ is entirely ordinary,
then $\W s u$. The source and intermediate term
may contain arbitrary marks.
\end{theorem}
\begin{proof}
We first derive $\EW s u$. Induct lexicographically
on the size of the right derivation and then that
of the left derivation. Sizes include all five
mutual judgments and every whole-source premise.
Thus a weak premise contained in a right strict
or multiset witness has smaller first coordinate.
When that coordinate decreases, the newly
constructed left proof may have any finite size.

If the left root is \textsc{Put}, either selection
rule, or an extra-list macro, apply the induction
hypothesis to its smaller continuation with the
right derivation fixed, and lift the result by
the same rule. After these cases have been
removed, inspect the right root.

\paragraph{Var, Abs, and Fun.}
For \textsc{Var}, $t=u$, so keep the left proof.
For \textsc{Abs}, the remaining left root is
\textsc{Abs} or \textsc{FAbs}. Cut their opened
body premises using the smaller right body
derivation and reapply that root.
For \textsc{Fun}, a remaining left
\textsc{Fun} composes the two occurrence
matchings by smaller-right cuts. A left
\textsc{Copy} keeps its strict head comparison,
which respects head equivalence, and transports
each whole-source premise by those cuts.
A left \textsc{Stat} transports both its strict
multiset and its whole-source premises through
the right weak matching. The precise multiset
transport is given below.

\paragraph{Put.}
The marking lemma gives
$\EW{\headmark(s)}{\headmark(t)}$.
The right premise compares $\headmark(t)$ with $u$
and is smaller than its \textsc{Put} conclusion.
Cut, and then remove the left total mark by
Lemma~\ref{lem:extended-marking}.

\paragraph{Copy and Stat.}
The intermediate term is marked, so the remaining
left root is \textsc{In}, \textsc{Marked1},
or \textsc{Marked2}. Each such comparison may
be retyped at the root without altering its
premises. Cut it against each smaller right
whole-source premise. Transitivity of head
precedence, and its compatibility with head
equivalence, give \textsc{Copy} when the right
rule is \textsc{Copy}.

If the right rule is \textsc{Stat}, a left
\textsc{Marked1} likewise yields \textsc{Copy}.
A left \textsc{In} transports its weak base
matching into the right strict multiset.
A left \textsc{Marked2} composes its strict
multiset with the right one. The whole-source
premises have already been transported;
\textsc{Stat} then gives the conclusion.

\paragraph{FAbs and target macros.}
A remaining \textsc{In}, \textsc{Marked1}, or
\textsc{Marked2} comparison survives opening
both owners with the same fresh extra variable.
In \textsc{In} append its reflexive comparison.
In the two precedence rules keep old
whole-source premises by \textsc{Drop} and
select the fresh extra for its own premise.
Cut with the smaller right body derivation,
then use \textsc{FAbs}.
For a right \textsc{Drop}, \textsc{Perm},
or \textsc{Copies}, transform the target of
the left comparison by
Lemma~\ref{lem:target-macros} and cut with
the smaller right premise.
The right root cannot be \textsc{In},
\textsc{Marked1}, or \textsc{Marked2},
because $u$ is ordinary.

\paragraph{Transport of strict multiset occurrences.}
Each final target occurrence is an ordinary
subterm of $u$. Consider an element path
$\EW a b$ followed by $\ES b c$.
Write $\mk(b)=b'$. Markability reflection
supplies $\mk(a)=a'$, and marking transport
gives $\EW{a'}{b'}$. Its cut against the
smaller right weak witness $\EW{b'}c$ gives
$\ES a c$. For a strict-then-weak path,
retain the original successful mark of $a$
and cut its weak witness. For a
strict-then-strict path, marking transports
$\EW{\mk(a)}b$ to
$\EW{\mk(a)}{\mk(b)}$, since the first
marked term is fixed by total marking.
The smaller right strict witness then
gives the composite strict edge.

To compose the occurrence partitions, keep
exactly those originally kept occurrences
whose intermediate images are kept again.
Remove every other source occurrence.
This removed set contains the first
removed set and is nonempty when the
first multiset is strict. Each target
occurrence is traced backwards through
the two supplier maps. A path that is
not the unique path between kept
occurrences has a strict edge, so the
element transport just proved supplies
its strict witness. If only the second
factor is strict, its nonempty removed
set has a nonempty inverse image under
the first weak bijection. The
strict-then-weak case is the same
construction with the second removed
set empty. Unused removed occurrences
are allowed. This verifies all strict,
weak--strict, and strict--weak cases,
including deletion without insertion.

\paragraph{The selection continuation.}
Let the right rule select an actual
argument $q$ of a marked owner $Q$,
apply it to prefix copies $Q_j$, obtain
$r$, and continue by $\EW r u$.
This continuation has smaller right
size. Let $P$ be the source of the
remaining left comparison with $Q$.

If that comparison is \textsc{In},
the selected base or extra has an
actual matching occurrence $p$ of $P$
with $\EW p q$. For each target prefix
$Q_j$, \textsc{In} compares the
corresponding source prefix with it,
and \textsc{Drop} restores the
remaining source extras. A whole $P$
copy at the required type therefore
dominates $Q_j$.
Lemma~\ref{lem:application-covariance}
gives a successful full call on $p$
with whole $P$ copies and compares
its result with $r$. Cut with the
smaller continuation $\EW r u$,
then use the original \textsc{Select}
rule at the actual argument $p$.

If the left comparison is
\textsc{Marked1} or \textsc{Marked2},
its whole-source premises give
$\EW{P_{\operatorname{type}(q)}}q$.
Here the type is the \emph{full
function type} of $q$.
The same root clause with truncated
target extras compares a whole $P$
with every requested $Q_j$.
Application covariance now gives
\[
 \EW{\ap{P_{\operatorname{type}(q)}}
                 {\text{whole }P\text{ copies}}}{r}.
\]
For an empty tuple the left result is
$P$ itself; otherwise it is exactly
$P$ with those copies appended.
Cut with $\EW r u$ and discharge
the appended full-prefix copies by
\textsc{Copies}. This finishes the
selection case.

Every root has been covered, and
every invocation used the stated
lexicographic decrease. We have
obtained $\EW s u$.
Apply Lemma~\ref{lem:conservative}
with the reflexive refinement of $s$
to obtain the required original
judgment $\W s u$.
\end{proof}
